\centerline{\bf \Huge Начало анализа}

\problem{1}
Пусть $\left(a_n\right)_{n=1}^{\infty}$ и $\left(b_n\right)_{n=1}^{\infty}$ - две сходящиеся к $A$ и $B$ соответственно последовательности. Тогда
(a) последовательность $\left(a_n+b_n\right)_{n=1}^{\infty}$ сходится к $A+B$; другими словами

$$
\lim _{n \rightarrow \infty}\left(a_n+b_n\right)=\lim _{n \rightarrow \infty} a_n+\lim _{n \rightarrow \infty} b_n .
$$

(b) последовательность $\left(a_n b_n\right)_{n=1}^{\infty}$ сходится к $A B$; другими словами

$$
\lim _{n \rightarrow \infty}\left(a_n b_n\right)=\lim _{n \rightarrow \infty} a_n \cdot \lim _{n \rightarrow \infty} b_n .
$$

(c) пусть $B \neq 0$, а также $b_n \neq 0$ при $n \geqslant m$. Тогда последовательность $\left(a_n / b_n\right)_{n=m}^{\infty}$ сходится к $A / B$; другими словами

$$
\lim _{n \rightarrow \infty} \frac{a_n}{b_n}=\frac{\lim _{n \rightarrow \infty} a_n}{\lim _{n \rightarrow \infty} b_n}
$$

\problem{2}
\textbf{Теорема о двух милиционерах.}
Последовательности $a_n, b_n$ и $c_n$ таковы, что $a_n \leqslant b_n \leqslant c_n$. Оказалось, что $a_n$ и $c_n$ стремятся к одному и тому же числу $d$. Докажите, что $b_n$ стремится к $d$. Сформулируйте и докажите аналогичное утверждение для функций. 

\problem{3}
(a) Пусть $f:[0,1] \rightarrow[0,1]$ непрерывная функция. Докажите, что найдется $x \in[0,1]$ такой, что $f(x)=x$.

(b) Пусть $f:[0,1] \rightarrow \mathbb{R}$ непрерывная функция, принимающая все значения из отрезка $[0,1]$. Докажите, что найдется $x \in[0,1]$ такой, что $f(x)=x$.

\problem{4}
Пусть $f:[0,1] \rightarrow[0,1]$ - непрерывная функция, причём $f(0)=f(1)=0$.

(a) Верно ли, что на графике $f$ найдётся горизонтальная хорда длины $1 / 3$ ?

(b) При каких $\ell$ обязательно найдётся горизонтальная хорда длины $\ell$ ?

\problem{5}
Пусть $m>k$ и $\mathrm{P}(\mathrm{x})=\mathrm{a}_{\mathrm{m}} \mathrm{x}^{\mathrm{m}}+\mathrm{a}_{\mathrm{m}-1} \mathrm{x}^{\mathrm{m}-1}+\ldots+\mathrm{a}_0$ и $\mathrm{Q}(\mathrm{x})=\mathrm{b}_{\mathrm{k}} \mathrm{x}^{\mathrm{k}}+\mathrm{b}_{\mathrm{k}-1} \mathrm{x}^{\mathrm{k}-1}+\ldots+\mathrm{b}_0-$ два многочлена (при этом $a_m, b_k \neq 0$ ). Найдите $\lim _{n \rightarrow \infty} \frac{P(n)}{Q(n)}$.

\problem{6}
С помощью критерия Коши проверьте сходимость последовательностей:

(a) $a_n=q^n$ при $|q|<1$

(b) $a_n=(-1)^n$